% ============================================================================
% Informe — Metodología de anonimización
% Oferta académica USTA 2021-1 a 2026-1 · Consultoría USTA 2026-II
% Universidad Santo Tomás — Documento consolidado: 1 de octubre de 2026
% ----------------------------------------------------------------------------
% Fuente del contenido: anonimizacion/METODOLOGIA_ANONIMIZACION.md
% (si cambia el Markdown, actualizar este archivo).
% A pedido del consultor (2026-10-01), esta versión OMITE las referencias a
% iCloud y al registro de la sesión de IA que sí están en el Markdown. Por eso
% su numeración es propia: P1-P5 (decisiones) y R1-R6 (hallazgos) no coinciden
% con P1-P6 y R1-R8 del Markdown. Al transcribir cambios, mantener la omisión.
% Formato tomado del Informe SAVIALAB (dos columnas, portada y anexos a una).
%
% Compilar dos veces (la segunda resuelve índice, referencias y la portada):
%     pdflatex Informe_Anonimizacion_USTA.tex
%     pdflatex Informe_Anonimizacion_USTA.tex
%
% No contiene datos personales: solo conteos y un ejemplo ficticio.
% Vive fuera de anonimizacion/ a propósito: 06_auditoria_final.py barre en
% bruto todo lo que hay ahí, y la tabla interna de un PDF tiene números de
% 10 dígitos que la auditoría tomaría por documentos de identidad.
% ============================================================================
\documentclass[11pt,a4paper,twoside,twocolumn]{article}

% --- Paquetes base ---
\usepackage[utf8]{inputenc}
\usepackage[T1]{fontenc}
\usepackage[spanish,es-tabla,es-noshorthands,es-nodecimaldot]{babel}
\usepackage{lmodern}
\usepackage{microtype}
\usepackage[a4paper,top=2.3cm,bottom=2.3cm,inner=2.0cm,outer=2.0cm,columnsep=0.7cm]{geometry}

\usepackage{graphicx}
\usepackage{booktabs}
\usepackage{array}
\usepackage{tabularx}
\usepackage{xltabular}
\usepackage{colortbl}
\usepackage{caption}
\usepackage{enumitem}
\usepackage{fancyhdr}
\setlength{\headheight}{14pt}
\usepackage{titlesec}
\usepackage[titles]{tocloft}
\setlength{\cftbeforesecskip}{3pt}
\setlength{\cftbeforesubsecskip}{0pt}
\usepackage{amsmath,amssymb}
\usepackage{xcolor}
\usepackage{listings}
\usepackage{fontawesome5}
\usepackage{tikz}
\usetikzlibrary{positioning,fit,arrows.meta,calc,backgrounds}
\usepackage[most]{tcolorbox}
\usepackage[hyphens]{url}
\usepackage{hyperref}

% --- Colores institucionales USTA ---
% Azul y dorado: los mismos de la plantilla USTA del proyecto
% (documentacion_repos_ustadistica/informes/_base/Informe_repo_base.tex).
% Rojo: el del escudo; aquí solo marca riesgos altos.
\definecolor{ustaAzul}{HTML}{003876}
\definecolor{ustaAzulOscuro}{HTML}{021952}
\definecolor{ustaDorado}{HTML}{FDB913}
\definecolor{ustaRojo}{HTML}{A6272E}
\definecolor{grisTexto}{HTML}{555555}
\colorlet{ustaDoradoTexto}{ustaDorado!70!black}

\hypersetup{
    colorlinks=true,
    linkcolor=ustaAzul,
    citecolor=ustaAzul,
    urlcolor=ustaAzul,
    pdftitle={Metodolog\'ia de anonimizaci\'on - Oferta acad\'emica USTA 2021-1 a 2026-1},
    pdfauthor={Javier Mauricio Sierra - Consultor\'ia USTA 2026-II},
    pdfsubject={Anonimizaci\'on de OFERTA\_ACADEMICA\_20211\_20261.xlsx}
}

% --- Rutas y nombres de archivo en monoespaciada, cortables en / _ . - ---
\urlstyle{tt}
\makeatletter
\g@addto@macro{\UrlBreaks}{\do\_\do\-\do\.}
\makeatother
\newcommand{\cod}[1]{\nolinkurl{#1}}

% --- Encabezados y pies doble cara ---
\pagestyle{fancy}
\fancyhf{}
\fancyhead[LE]{\small\textit{Metodología de anonimización --- Oferta académica USTA}}
\fancyhead[RO]{\small\textit{\leftmark}}
\fancyfoot[LE,RO]{\textcolor{ustaAzul}{\thepage}}
\fancyfoot[RE,LO]{\small\textcolor{ustaAzul}{Universidad Santo Tomás} \textcolor{ustaDoradoTexto}{$\cdot$} \textcolor{grisTexto}{Consultoría 2026-II}}
\renewcommand{\headrulewidth}{0.6pt}
\renewcommand{\footrulewidth}{0.4pt}
\renewcommand{\headrule}{{\color{ustaDorado}\hrule width\headwidth height\headrulewidth\vskip-\headrulewidth}}
\renewcommand{\footrule}{{\color{ustaAzul}\vskip-\footruleskip\vskip-\footrulewidth\hrule width\headwidth height\footrulewidth\vskip\footruleskip}}

% --- Espaciado y títulos ---
\setlength{\parskip}{0.5\baselineskip}
\setlength{\parindent}{0pt}
\emergencystretch=3em
\hyphenpenalty=200
\tolerance=1000
\raggedbottom
\titleformat{\section}{\normalfont\large\bfseries\raggedright\color{ustaAzul}}{\thesection.}{0.5em}{}[{\color{ustaDorado}\titlerule[0.8pt]}]
\titleformat{\subsection}{\normalfont\normalsize\bfseries\raggedright\color{ustaAzulOscuro}}{\thesubsection.}{0.5em}{}
\titleformat{\subsubsection}{\normalfont\small\bfseries\raggedright\color{ustaAzulOscuro}}{\thesubsubsection.}{0.5em}{}
\titlespacing*{\subsubsection}{0pt}{1.2ex plus .5ex}{0.4ex}

% Pasos del proceso: rótulo azul + título
\newcommand{\paso}[2]{\subsubsection*{\colorbox{ustaAzul}{\textcolor{white}{\scriptsize\sffamily\bfseries PASO #1}}\enspace #2}}

% --- Flotantes a dos columnas ---
\renewcommand{\topfraction}{0.9}
\renewcommand{\dbltopfraction}{0.72}
\renewcommand{\textfraction}{0.07}
\renewcommand{\floatpagefraction}{0.8}
\renewcommand{\dblfloatpagefraction}{0.6}
\setcounter{topnumber}{3}
\setcounter{dbltopnumber}{3}
\setcounter{totalnumber}{5}

% --- Tablas ---
\captionsetup{font=small,labelfont={bf,color=ustaAzul},labelsep=period}
\arrayrulecolor{ustaAzul}
\renewcommand{\arraystretch}{1.18}
\newcolumntype{L}{>{\raggedright\arraybackslash}X}
\newcolumntype{P}[1]{>{\raggedright\arraybackslash}p{#1}}
\newcolumntype{Q}[1]{>{\raggedleft\arraybackslash}p{#1}}
\newcommand{\cab}[1]{\textcolor{ustaAzul}{\textbf{#1}}}

% --- Marcas de nivel de riesgo (sustituyen los emojis del Markdown) ---
\newcommand{\punto}[1]{\tikz[baseline=-0.55ex]\fill[#1] (0,0) circle (0.55ex);}
\newcommand{\alto}{\punto{ustaRojo}}
\newcommand{\medio}{\punto{ustaDorado}}
\newcommand{\nose}{\textcolor{ustaRojo}{\faBan}\,\textbf{No}}

% --- Cajas ---
\tcbset{
  cajabase/.style={enhanced,boxrule=0pt,arc=0pt,outer arc=0pt,
    left=6pt,right=6pt,top=4pt,bottom=4pt,before skip=6pt,after skip=8pt,
    fontupper=\small}
}
\newtcolorbox{nota}{cajabase,colback=ustaAzul!5,colframe=ustaAzul,leftrule=2.5pt}
\newtcolorbox{alerta}{cajabase,colback=ustaRojo!5,colframe=ustaRojo,leftrule=2.5pt}
\newtcolorbox{dorada}{cajabase,colback=ustaDorado!12,colframe=ustaDorado,leftrule=2.5pt}

% --- Código ---
\lstset{
  basicstyle=\ttfamily\scriptsize,
  columns=fullflexible,
  keepspaces=true,
  breaklines=true,
  breakatwhitespace=true,
  showstringspaces=false,
  upquote=true,
  commentstyle=\color{grisTexto}\itshape,
  morecomment=[l]{\#},
  literate={á}{{\'a}}1 {é}{{\'e}}1 {í}{{\'i}}1 {ó}{{\'o}}1 {ú}{{\'u}}1
           {Í}{{\'I}}1 {ñ}{{\~n}}1
}
\newtcblisting{codigo}[1][]{enhanced,boxrule=0pt,arc=0pt,outer arc=0pt,
  colback=ustaAzul!4,colframe=ustaAzul,leftrule=2.5pt,boxsep=0pt,left=4pt,right=2pt,top=3pt,bottom=3pt,
  listing only,listing options={basicstyle=\ttfamily\scriptsize,columns=fullflexible,keepspaces=true,breaklines=true,#1}}

% ============================================================================
\begin{document}
% ============================================================================

% --- Portada a página completa ---
\onecolumn
\begin{titlepage}
\begin{tikzpicture}[remember picture,overlay]
  \fill[ustaAzulOscuro] (current page.north west) rectangle ([yshift=-3.1cm]current page.north east);
  \fill[ustaDorado] ([yshift=-3.1cm]current page.north west) rectangle ([yshift=-3.28cm]current page.north east);
  \node[text=white,font=\LARGE\bfseries] at ([yshift=-1.3cm]current page.north) {UNIVERSIDAD SANTO TOMÁS};
  \node[text=ustaDorado,font=\normalsize] at ([yshift=-2.15cm]current page.north)
       {Facultad de Estadística \enspace$\cdot$\enspace Asignatura Consultoría \enspace$\cdot$\enspace 2026-II};
  \fill[ustaAzulOscuro] (current page.south west) rectangle ([yshift=1.0cm]current page.south east);
  \fill[ustaDorado] ([yshift=1.0cm]current page.south west) rectangle ([yshift=1.12cm]current page.south east);
  \node[text=white,font=\footnotesize] at ([yshift=0.5cm]current page.south) {Documento de uso interno \enspace$\cdot$\enspace no contiene datos personales};
\end{tikzpicture}
\centering
\vspace*{2.6cm}

{\large\color{ustaAzul} Proyecto de consultoría USTA 2026-II}\\[0.9cm]

{\color{ustaDorado}\rule{\textwidth}{1.2pt}}\\[0.35cm]
{\LARGE\bfseries\color{ustaAzulOscuro} Metodología de anonimización}\\[0.2cm]
{\LARGE\bfseries\color{ustaAzulOscuro} Oferta académica USTA 2021-1 a 2026-1}\\[0.35cm]
{\large\itshape Qué se hizo, por qué, cómo se comprobó\\ y qué no garantiza}\\[0.3cm]
{\color{ustaDorado}\rule{\textwidth}{1.2pt}}\\[1.0cm]

{\large\bfseries\color{ustaAzul} Responsable}\\[0.15cm]
{\Large Javier Mauricio Sierra}\\[0.8cm]

{\large\bfseries\color{ustaAzul} Fechas}\\[0.15cm]
{Ejecución: 2 y 3 de septiembre de 2026 \quad$\cdot$\quad Documento consolidado: 1 de octubre de 2026}\\[0.9cm]

\begin{tabular}{rl}
\textbf{Archivo de origen:} & \cod{OFERTA_ACADEMICA_20211_20261.xlsx} \\
\textbf{SHA-256 del original:} & \texttt{f447e0d2\ldots a9d9} (sin cambios) \\
\textbf{Propósito:} & Simular el cobro del pregrado por créditos \\
\textbf{Universo del estudio:} & 60 programas de pregrado $\cdot$ 18.563 estudiantes \\
\textbf{Entregables:} & \cod{matriculas_anon} (78.625 $\times$ 20, principal) \\
                     & \cod{inscripciones_anon} (448.489 $\times$ 14, detalle) \\
\textbf{Método:} & HMAC-SHA256 con llave efímera destruida; supresión, \\
                 & generalización, barajado y medición del riesgo \\
\textbf{Estado:} & Ejecutado; cinco decisiones pendientes (P1--P5) \\
\end{tabular}

\vfill
{\footnotesize Este documento es la puerta de entrada a la anonimización del archivo de oferta
académica entregado por la Universidad Santo Tomás. Los detalles están en la bitácora
(\cod{EVIDENCIA.md}), el informe de riesgo (\cod{RIESGO_REIDENTIFICACION.md}), el diccionario de
datos (\cod{entrega/DICCIONARIO_DATOS.md}) y los scripts (\cod{anonimizacion/scripts/}). Todas las
cifras se recalcularon desde los datos el 1 de octubre de 2026; las que no salen de otro script
las reproduce \cod{08_revision_documental.py}.}\\[1.4cm]
\end{titlepage}

% --- Índice y guía de lectura a una columna ---
\begingroup
\setlength{\parskip}{0pt}
\linespread{0.94}\selectfont
\tableofcontents
\endgroup

\newpage

% --- Resto del documento a dos columnas ---
\twocolumn

\begin{nota}
{\normalsize\bfseries\color{ustaAzul} Cómo leer este documento}\par\smallskip
\footnotesize
\setlength{\tabcolsep}{3pt}
\begin{tabularx}{\linewidth}{@{}L P{2.3cm} r@{}}
\toprule
\cab{Si usted es\ldots} & \cab{Lea} & \cab{Tiempo} \\
\midrule
Directivo/a o responsable de protección de datos de la USTA & §\ref{sec:resumen}, §\ref{sec:riesgo}, §\ref{sec:limites} y §\ref{sec:custodia} & 10 min \\
Analista que va a usar los datos & §\ref{sec:resumen}, §\ref{sec:columnas}, §\ref{sec:riesgo} y el diccionario & 15 min \\
Auditor/a técnico/a & Todo; después la bitácora y los scripts & 1 h \\
Quien retome el proyecto & §\ref{sec:custodia}, §\ref{sec:verificar} y Anexo~\ref{app:mapa} & 15 min \\
\bottomrule
\end{tabularx}
\par\smallskip
Los términos técnicos están en el glosario (Anexo~\ref{app:glosario}).
\end{nota}

% ============================================================================
\section{Resumen ejecutivo}\label{sec:resumen}
% ============================================================================

\textbf{Para qué.} El estudio quiere estimar qué habría pasado si la USTA hubiera cobrado el
pregrado \textbf{por créditos} en lugar de con su esquema actual de \textbf{media matrícula y
matrícula completa}. Para eso hacen falta, por cada estudiante y semestre, los créditos
inscritos y el valor efectivamente pagado. No hace falta saber quién es nadie.

\textbf{Qué se recibió.} Un Excel de 656.737 filas ---una por cada materia que inscribió cada
estudiante--- y 52 columnas, con 11 periodos (2021-1 a 2026-1) y datos de 41.589 personas:
nombres, cédulas, correos, celulares, notas, estado académico y valor pagado.

\textbf{Qué se entrega.} Dos tablas anonimizadas de los \textbf{18.563 estudiantes de los 60
programas de pregrado} (Tabla~\ref{tab:entregables}).

\begin{table}[!ht]
\centering
\caption{Entregables de la anonimización (CSV comprimido, \texttt{.csv.gz}).}
\label{tab:entregables}
\footnotesize
\setlength{\tabcolsep}{3.5pt}
\begin{tabularx}{\columnwidth}{@{}P{3.3cm} L Q{1.45cm}@{}}
\toprule
\cab{Tabla y uso} & \cab{Una fila por\ldots} & \cab{Filas $\times$ col.} \\
\midrule
\cod{matriculas_anon}\newline \textbf{Principal.} Basta para la simulación & estudiante $\times$ programa $\times$ periodo & 78.625 $\times$ 20 \\
\cod{inscripciones_anon}\newline Detalle. Solo con justificación & materia inscrita & 448.489 $\times$ 14 \\
\bottomrule
\end{tabularx}
\end{table}

\textbf{Cómo, en cinco ideas.}

\begin{enumerate}[leftmargin=*,itemsep=2pt,topsep=0pt]
  \item \textbf{Minimizar según el propósito.} Lo que no sirve para costear no sale: se eliminaron
  31 de las 52 columnas ---los 6 identificadores, las 19 de notas, el grupo, la sede y 4 que se
  sustituyeron por versiones recodificadas--- y todo lo que no es pregrado.
  \item \textbf{Seudónimos con llave destruida.} La cédula y el código de estudiante se
  reemplazaron por códigos (\texttt{E…}, \texttt{M…}) calculados con HMAC-SHA256 y una llave aleatoria
  de 256 bits que solo existió en la memoria del proceso y se destruyó al terminar. A partir del
  entregable, nadie ---ni siquiera quien hizo el proceso--- puede convertir un código en una cédula.
  \item \textbf{Generalizar lo que singulariza.} Valor pagado redondeado al millar, ingresos
  anteriores a 2015 agrupados, reintegros topados, subsidio y sanción disciplinaria reducidos a
  categorías.
  \item \textbf{Barajar las filas}, para que la posición no delate el orden del archivo original.
  \item \textbf{Medir el riesgo que queda} en lugar de suponerlo, y documentarlo.
\end{enumerate}

\textbf{Qué está comprobado.} El original nunca se modificó (mismo SHA-256 al inicio, al final y
hoy). La extracción es fiel (1.300 de 1.300 celdas). Los entregables no contienen el nombre de
ninguno de los 18.563 estudiantes, ni correos, cédulas o celulares. El dinero y los créditos
cuadran contra el original: en cada uno de los 497 grupos de programa y periodo, los valores
entregados son exactamente los originales redondeados al millar (error $\leq$ \$500 por
matrícula), y los 1.158.935 créditos son idénticos.

\begin{alerta}
\textbf{\color{ustaRojo}\faExclamationTriangle\ Qué no está garantizado} (detalle en
§\ref{sec:riesgo} y §\ref{sec:limites}).
\begin{itemize}[leftmargin=*,itemsep=2pt,topsep=2pt]
  \item \textbf{Riesgo residual.} En el 2,91\,\% de las matrículas, la combinación de periodo,
  programa y nivel la comparten menos de 5 matrículas (contando la propia), y el 71\,\% de los
  estudiantes tiene un conjunto de materias único. Es el precio de conservar el detalle que el
  estudio necesita.
  \item \textbf{Frente a quien tenga el archivo fuente, no es anónimo.} La USTA, o este equipo
  mientras conserve el original, puede volver a ligar los registros con las personas.
\end{itemize}
\end{alerta}

\textbf{Decisiones pendientes} (§\ref{sec:pendientes}). Del consultor: rebarajar los entregables
antes de que salgan del equipo y destruir los intermedios en claro. De la USTA: autorizar el
tratamiento, aceptar por escrito el riesgo residual y decidir el destino del original.

% ============================================================================
\section{Contexto y marco}\label{sec:contexto}
% ============================================================================

\subsection{El propósito manda}

La regla que organizó todo el trabajo: \textbf{cada columna tiene que justificarse por el
análisis tarifario, o sale.} Es la medida de privacidad más fuerte disponible, porque lo que no
se entrega no se puede filtrar, y además no cuesta nada al análisis (Tabla~\ref{tab:necesita}).

\begin{table}[!ht]
\centering
\caption{Qué necesita la simulación y qué no.}
\label{tab:necesita}
\footnotesize
\begin{tabularx}{\columnwidth}{@{}L L@{}}
\toprule
\cab{La simulación necesita} & \cab{No necesita} \\
\midrule
Créditos inscritos por matrícula & Nombre, cédula, correos, celular \\
Valor pagado por matrícula & Notas, promedios, fallas \\
Programa, facultad, periodo, nivel & Grupo de clase \\
Estado académico, reintegro, subsidio (como controles) & Sede (es una sola) \\
Seguir al mismo estudiante entre semestres & Saber quién es ese estudiante \\
\bottomrule
\end{tabularx}
\end{table}

La última fila es la tensión central: el estudio necesita reconocer que dos matrículas son de la
misma persona, sin saber de quién. Eso es lo que resuelve la seudonimización
(§\ref{sec:proceso}, paso~4).

\subsection{Anonimizar y seudonimizar no son lo mismo}\label{sec:anon-seud}

\begin{itemize}[leftmargin=*,itemsep=2pt,topsep=0pt]
  \item \textbf{Seudonimizar} es reemplazar los identificadores por códigos. Si alguien puede
  reconstruir el vínculo ---con una llave, una tabla de correspondencias o el archivo original---
  el dato sigue siendo personal.
  \item \textbf{Anonimizar} es que nadie pueda reconstruirlo con medios razonables.
\end{itemize}

Este proceso destruyó la llave: quien solo recibe los entregables no tiene forma de convertir un
código en una cédula. Lo que sí conserva es la posibilidad de \textbf{reconocer} a alguien por
sus atributos ---programa, semestre, materias, valor pagado--- si ya sabe cosas de esa persona;
ese riesgo se mide en la §\ref{sec:riesgo}. Quien conserva el archivo fuente, en cambio, puede
volver a ligar los registros (§\ref{sec:fuente}): para él el resultado es
\textbf{seudonimizado}. Ninguno de los dos casos es un anonimato sin matices, y la distinción
importa para la calificación jurídica del entregable.

\subsection{Marco normativo de referencia}

Esta sección orienta; \textbf{no sustituye un concepto jurídico} de la USTA.

\begin{itemize}[leftmargin=*,itemsep=2pt,topsep=0pt]
  \item \textbf{Ley 1581 de 2012} y su reglamentación (Decreto 1377 de 2013, compilado en el
  Decreto 1074 de 2015). Define dato personal como información vinculada o que pueda asociarse
  a personas determinadas o \emph{determinables}. Un dato que ya no puede asociarse a nadie deja
  de serlo.
  \item \textbf{Roles.} La USTA es responsable del tratamiento; el consultor actúa como
  encargado.
  \item \textbf{Principios aplicados} (art.~4): \emph{finalidad} (solo lo necesario para el
  estudio), \emph{acceso y circulación restringida} (los intermedios no circulan),
  \emph{seguridad} y \emph{confidencialidad}.
  \item \textbf{Mientras existan los intermedios en claro hay tratamiento de datos personales},
  con la finalidad acotada a este ejercicio.
  \item \textbf{Pendiente desde el inicio:} confirmar que la USTA autorizó tratar este archivo con
  fines de consultoría (§\ref{sec:pendientes}, P3).
\end{itemize}

% ============================================================================
\section{El archivo de origen}\label{sec:origen}
% ============================================================================

\subsection{Estructura}

La Tabla~\ref{tab:estructura} resume la forma del archivo recibido.

\begin{table}[!ht]
\centering
\caption{Estructura del archivo de origen.}
\label{tab:estructura}
\footnotesize
\begin{tabularx}{\columnwidth}{@{}P{2.25cm} L@{}}
\toprule
\cab{Atributo} & \cab{Valor} \\
\midrule
Archivo & \cod{OFERTA_ACADEMICA_20211_20261.xlsx}, 142.819.604 bytes \\
Hojas & 1 (\texttt{Exportar Hoja de Trabajo}) \\
Filas de datos & 656.737 --- una por materia inscrita por un estudiante en un periodo \\
Columnas & 52 (\cod{A}\ldots\cod{AZ}) \\
Periodos & 11, de \cod{20211} a \cod{20261} \\
Personas & 41.592 cédulas distintas tal como vienen; \textbf{41.589} tras normalizar (§\ref{sec:proceso}, paso~4) \\
Nombres distintos & 41.571: hay homónimos, por eso hay menos nombres que personas \\
Códigos de estudiante & 54.637; 9.438 personas tienen más de uno (cursan o cursaron varios programas) \\
Programas & 576, de pregrado, posgrado y educación continuada \\
Tamaño real & 1,18 GB de XML descomprimido en una sola hoja; 266.647 cadenas compartidas \\
\bottomrule
\end{tabularx}
\end{table}

\subsection{Qué contiene, por nivel de riesgo}

No todas las columnas pesan igual. La Tabla~\ref{tab:niveles} las agrupa por el daño que harían si se filtraran.

\begin{table}[!ht]
\centering
\caption{Columnas del original por nivel de riesgo.}
\label{tab:niveles}
\footnotesize
\setlength{\tabcolsep}{3.5pt}
\begin{tabularx}{\columnwidth}{@{}P{1.9cm} L P{2.2cm}@{}}
\toprule
\cab{Tipo} & \cab{Columnas} & \cab{Riesgo} \\
\midrule
Identifi\-cadores directos & nombres, cédula, correo institucional, correo personal, celular, código & \alto\ Identifican por sí solos \\
Datos de alto impacto & 19 de notas, promedios y fallas; estado (incluye sanción disciplinaria); Ser Pilo Paga & \alto\ Dañinos si se filtran \\
Cuasi-identi\-ficadores & programa, periodo, nivel, periodo de ingreso, grupo, valor pagado, facultad & \medio\ Identifican combinados \\
Carga analítica & créditos, materias, pensum, valor pagado & Necesarios para el estudio \\
\bottomrule
\end{tabularx}
\end{table}

\subsection{Hallazgos sobre el archivo que condicionaron el diseño}

\begin{enumerate}[leftmargin=*,itemsep=2pt,topsep=0pt]
  \item \textbf{No se puede editar en sitio.} 1,18 GB de XML en una hoja agota la memoria de las
  herramientas habituales. Y aunque se pudiera, borrar celdas de un Excel \textbf{no borra los
  textos}: quedan en \cod{xl/sharedStrings.xml} como residuo recuperable. Por eso el flujo es
  \emph{extraer $\to$ anonimizar $\to$ escribir un archivo nuevo}, nunca \emph{editar}.
  \item \textbf{Metadatos con origen.} \cod{xl/workbook.xml} guarda la ruta de la cuenta de
  Registro y Control desde la que se exportó. No pasa a los CSV, que no llevan metadatos de
  Office.
  \item \textbf{La persona es la cédula, no el código.} Quien cambia o suma programas recibe otro
  código; seudonimizar por código partiría la historia de 9.438 personas.
  \item \textbf{Cédulas duplicadas por ceros a la izquierda.} Tres personas aparecían dos veces
  (con y sin ceros iniciales). De ahí las 41.592 «personas» que en realidad son 41.589.
  \item \textbf{Una columna mal rotulada.} La columna \cod{20221} no es un periodo ni está vacía:
  tiene 283.443 calificaciones.
  \item \textbf{\cod{VAL_PAGADO} se repite en cada materia} del mismo estudiante y periodo: es el
  valor de la matrícula del semestre, no de la materia.
  \item \textbf{\cod{PAGO} es constante en la práctica} («Con Pago» en el 99,999\,\%), así que no
  distingue media matrícula de matrícula completa; eso hay que derivarlo de \cod{VAL_PAGADO}.
  \item \textbf{La sede es una sola} (\cod{CAU} constante) y \textbf{\cod{PER_INGRESO} solo está
  informado en el 38,8\,\% de las filas} del archivo completo.
\end{enumerate}

% ============================================================================
\section{Principios y decisiones de diseño}\label{sec:decisiones}
% ============================================================================

La Tabla~\ref{tab:decisiones} recoge cada decisión con sus alternativas y quién la tomó.
«Consultor» indica una decisión explícita del responsable del proyecto; «técnica», una decisión
de ejecución de quien implementó, dentro del plan aprobado o como desviación documentada en la
§\ref{sec:desviaciones}.

\begin{table*}[tp]
\centering
\caption{Las doce decisiones de diseño.}
\label{tab:decisiones}
\footnotesize
\setlength{\tabcolsep}{4pt}
\begin{tabularx}{\textwidth}{@{}r P{3.9cm} P{3.4cm} L P{2.3cm}@{}}
\toprule
\cab{\#} & \cab{Decisión} & \cab{Alternativas descartadas} & \cab{Por qué} & \cab{Quién} \\
\midrule
1 & \textbf{El original no se toca}: solo lectura (\cod{444}), checksum registrado, todo sobre una copia & Trabajar sobre el original & Es la prueba de que los entregables vienen de él & Consultor \\
2 & Extraer $\to$ anonimizar $\to$ \textbf{escribir archivos nuevos} & Borrar columnas en el Excel & El residuo en \cod{sharedStrings.xml} y los metadatos & Técnica \\
3 & \textbf{Seudonimización irreversible}: llave destruida & Llave guardada en bóveda (reversible); hash simple; numeración consecutiva & Sin llave nadie puede revertirla; un hash simple se rompe en horas (§\ref{sec:proceso}, paso~4) & Consultor \\
4 & \textbf{La llave de sujeto es la cédula normalizada} & El código de estudiante & El código partiría la trayectoria de 9.438 personas & Técnica \\
5 & \textbf{Solo pregrado}: 60 programas & Todo el archivo & Es el alcance del estudio, y de paso elimina cientos de programas con 1 a 5 personas & Consultor, con 9 decisiones caso a caso \\
6 & \textbf{Programas pequeños con su nombre real} & Colapsarlos en \texttt{OTRO PREGRADO} & Conservar la desagregación por programa; riesgo aceptado a sabiendas por el consultor & Consultor \\
7 & \textbf{Suprimir las notas} (19 columnas) & Conservarlas & Irrelevantes para costear y de alto impacto & Técnica \\
8 & \textbf{Dos tablas}, la agregada como principal & Una sola tabla de detalle & El detalle por materia concentra el riesgo de trayectoria & Técnica \\
9 & \textbf{CSV comprimido}, \cod{;}, UTF-8 con BOM & Parquet o Excel & Lo pidió el consultor; abre directo en Excel en español & Consultor \\
10 & \textbf{Barajar con semilla fija} & No barajar; semilla secreta & Reproducibilidad. \emph{En retrospectiva fue un error}: la semilla publicada permite deshacer el barajado a quien tenga el fuente (§\ref{sec:fuente}) & Técnica \\
11 & \textbf{Redondear el valor al millar} & Bandas de valor & Error $\leq$ \$500 en matrículas de millones; no afecta el costeo & Técnica \\
12 & \textbf{Seudonimizar, filtrar y escribir en un solo proceso} & Un script por tarea, como decía el plan & Con llave efímera, cada ejecución da códigos distintos & Técnica \\
\bottomrule
\end{tabularx}
\end{table*}

% ============================================================================
\section{El proceso, paso a paso}\label{sec:proceso}
% ============================================================================

\begin{figure*}[tp]
\centering
\begin{tikzpicture}[
  font=\scriptsize,
  >={Stealth[length=2.2mm]},
  caja/.style={rectangle, rounded corners=2pt, align=center, inner sep=3pt, minimum height=1.0cm, text width=#1},
  claro/.style={caja=#1, draw=ustaRojo, fill=ustaRojo!7},
  anon/.style={caja=#1, draw=ustaAzul, very thick, fill=ustaAzul!9},
  util/.style={caja=#1, draw=ustaDoradoTexto, fill=ustaDorado!14},
  tarea/.style={rectangle, rounded corners=2pt, align=center, inner sep=3pt, text width=3.3cm, fill=white},
  etapa/.style={font=\footnotesize\bfseries, text=white, fill=ustaAzul, rounded corners=2pt, align=center, inner sep=3pt, minimum height=0.95cm, text width=#1},
  flujo/.style={->, thick, draw=ustaAzulOscuro},
  verif/.style={->, dashed, thick, draw=ustaAzul!70}
]
% Etapas (columnas en x = 0; 4,15; 8,55; 13,3)
\node[etapa=3.0cm] at (0,0)     {1 $\cdot$ Congelar\\ y copiar};
\node[etapa=3.3cm] at (4.15,0)  {2 $\cdot$ Extraer y\\ fijar el universo};
\node[etapa=3.6cm] at (8.55,0)  {3 $\cdot$ Anonimizar\\ en un solo proceso};
\node[etapa=3.8cm] at (13.3,0)  {4 $\cdot$ Verificar\\ y documentar};

% Etapa 1
\node[claro=3.0cm] (orig) at (0,-1.6) {\textbf{Original}\\ solo lectura (444)\\ \cod{OFERTA_ACADEMICA_20211_20261.xlsx}\\ SHA-256 registrado};
\node[claro=3.0cm] (copia) at (0,-3.9) {\textbf{Copia de trabajo}\\ \cod{trabajo/OFERTA_COPIA.xlsx}\\ idéntica: mismo checksum};
\draw[flujo] (orig) -- node[right,font=\scriptsize\ttfamily]{cp} (copia);

% Etapa 2
\node[util=3.3cm] (v01b) at (4.15,-1.6) {\cod{01b_verificar.py}\\ compara con el \textbf{original}:\\ 1.300 de 1.300 celdas iguales};
\node[claro=3.3cm] (raw) at (4.15,-3.9) {\textbf{Extracción en texto crudo}\\ \cod{_INSEGURO/oferta_raw.parquet}\\ \texttt{01\_extraer.py}, en streaming\\ 656.737 filas $\cdot$ 19,7 MB};
\node[util=3.3cm] (filtro) at (4.15,-6.35) {\cod{02_catalogo_programas.py}\\ \cod{02c_filtro_validado.py}\\ 576 programas $\to$ \textbf{60} de pregrado (Checkpoint A)};
\draw[flujo] (copia) -- (raw);
\draw[verif] (orig) -- (v01b);
\draw[verif] (raw) -- (v01b);
\draw[flujo] (raw) -- (filtro);

% Etapa 3
\node[font=\scriptsize\ttfamily\bfseries, text=ustaAzul] (t0) at (8.55,-0.95) {03\_anonimizar.py};
\node[tarea] (t4) at (8.55,-2.0) {\textbf{T4 $\cdot$ Seudonimizar}\\ HMAC-SHA256; llave de 256 bits solo en memoria, \textcolor{ustaRojo}{\textbf{destruida}} al terminar};
\node[tarea] (t5) at (8.55,-3.9) {\textbf{T5 $\cdot$ Filtrar, suprimir y generalizar}\\ salen 31 de 52 columnas};
\node[tarea] (t6) at (8.55,-5.65) {\textbf{T6 $\cdot$ Agregar, barajar y escribir}\\ CSV con \texttt{;}, \mbox{UTF-8} con BOM, gzip};
\draw[->, draw=ustaAzul!70] (t4) -- (t5);
\draw[->, draw=ustaAzul!70] (t5) -- (t6);
\begin{scope}[on background layer]
  \node[draw=ustaAzul, very thick, fill=ustaAzul!7, rounded corners=3pt, inner sep=5pt, fit=(t0)(t4)(t5)(t6)] (proc) {};
\end{scope}
\coordinate (yfiltro) at (0,-4.6);
\draw[flujo] (raw.east) -- (raw.east -| proc.west);
\draw[flujo] (filtro.east) -| (6.25,-4.6) -- (proc.west |- yfiltro);

% Etapa 4
\node[anon=3.8cm] (mat) at (13.3,-1.45) {\cod{matriculas_anon.csv.gz}\\ 78.625 $\times$ 20 $\cdot$ \textbf{principal}};
\node[anon=3.8cm] (ins) at (13.3,-2.85) {\cod{inscripciones_anon.csv.gz}\\ 448.489 $\times$ 14 $\cdot$ detalle};
\draw[flujo] (proc.east |- mat) -- (mat);
\draw[flujo] (proc.east |- ins) -- (ins);
\node[util=3.8cm, align=left] (chk) at (13.3,-5.35) {%
  \renewcommand{\arraystretch}{1.05}%
  \begin{tabular}{@{}>{\ttfamily}l@{\;}l@{}}
  03b    & Checkpoint B: 23/23 \\
  04     & informe de riesgo \\
  05     & diccionario, checksums \\
  06     & auditoría final \\
  07     & prefijos \texttt{E…} y \texttt{M…} \\
  05/06  & versión final \\
  08     & revisión 2026-10-01 \\
  \end{tabular}};
\draw[verif] (ins) -- (chk);

% Leyenda
\node[anchor=west] (l1) at (-1.6,-7.75) {\fcolorbox{ustaRojo}{ustaRojo!7}{\rule{0pt}{4pt}\hspace{9pt}}\ datos personales en claro};
\node[anchor=west] (l2) at ([xshift=2.5mm]l1.east) {\fcolorbox{ustaAzul}{ustaAzul!9}{\rule{0pt}{4pt}\hspace{9pt}}\ entregable anonimizado};
\node[anchor=west] (l3) at ([xshift=2.5mm]l2.east) {\fcolorbox{ustaDoradoTexto}{ustaDorado!14}{\rule{0pt}{4pt}\hspace{9pt}}\ verificación};
\draw[flujo] ([xshift=2.5mm]l3.east) -- ++(0.6,0) node[right] (l4) {flujo de datos};
\draw[verif] ([xshift=2.5mm]l4.east) -- ++(0.6,0) node[right] {comprobación};
\end{tikzpicture}
\caption{Flujo de la anonimización. El original solo se lee; todo el trabajo se hace sobre una
copia. La llave HMAC existe únicamente dentro de \texttt{03\_anonimizar.py}, que seudonimiza,
filtra y escribe en un solo proceso.}
\label{fig:flujo}
\end{figure*}

Cada paso tiene un objetivo, una acción y una comprobación que lo cierra; un paso no se daba por
terminado sin su comprobación (Figura~\ref{fig:flujo}). Entre paréntesis va la tarea
correspondiente del plan (\cod{PLAN_ANONIMIZACION.md}). Los \emph{checkpoints} A, B y C eran
puntos de control en los que el consultor revisaba los resultados y aprobaba antes de seguir.

\paso{0}{Perfilado (solo lectura)}

\textbf{Objetivo:} dimensionar el archivo antes de decidir nada. Se leyó el original en modo de
solo lectura, sin escribir nada. Produjo las cifras de la §\ref{sec:origen} y descubrió las
restricciones técnicas.

\paso{1}{Congelar el original y copiarlo (Tarea 1)}

\begin{itemize}[leftmargin=*,itemsep=2pt,topsep=0pt]
  \item SHA-256 del original registrado:\\
  {\scriptsize\texttt{f447e0d24e716b09a9abea1655dec6a0\\c2e143bd418fd93115e826c5a0eaa9d9}}.
  \item Permisos cambiados de \cod{644} a \cod{444}. Es la \textbf{única} acción aplicada al
  original: cambia un atributo del sistema de archivos, no el contenido, y el checksum posterior
  lo demuestra.
  \item Copia bit a bit en \cod{trabajo/OFERTA_COPIA.xlsx}, con checksum idéntico.
\end{itemize}

\textbf{Comprobación:} el mismo SHA-256 en el original al inicio, tras la extracción, en la
auditoría final y en la revisión del 2026-10-01.

\paso{2}{Extraer a un formato manejable (Tarea 3)}

\cod{01_extraer.py} lee la copia ---nunca el original--- en \emph{streaming}: recorre el XML fila
a fila y libera la memoria de cada una, escribiendo bloques de 50.000 filas a Parquet.
54 segundos; resultado de 19,7 MB.

Los valores se guardan \textbf{como texto, exactamente como están en el XML}. El plan pedía darles
tipo aquí, pero convertir a número antes de comparar habría hecho inexacta la verificación (un
\cod{4.4000000000000004} dejaría de coincidir consigo mismo). El tipado se hizo después.

\textbf{Comprobación} (\cod{01b_verificar.py}): 25 filas elegidas al azar (semilla 20260903) se
compararon celda a celda \textbf{contra el original}, no contra la copia; así la muestra pone a
prueba la cadena completa, original $\to$ copia $\to$ extracción. Resultado: \textbf{1.300 de
1.300 celdas idénticas}. Además, 656.737 filas y los conteos de no nulos de 14 columnas clave
coinciden con el perfilado. Que la copia es idéntica al original lo garantiza, aparte, su
checksum.

\paso{3}{Definir el universo del estudio (Tarea 2 y Checkpoint A)}

¿Qué es pregrado? El archivo no lo dice. \cod{02_catalogo_programas.py} clasificó los 576
programas con reglas explícitas y marcó lo dudoso para revisión humana en vez de adivinar:

\begin{itemize}[leftmargin=*,itemsep=1pt,topsep=0pt]
  \item \textbf{Posgrado:} el nombre dice especialización, maestría o doctorado.
  \item \textbf{Educación continuada:} diplomado, curso, taller, congreso, seminario y similares.
  \item \textbf{Movilidad:} intercambio, homologación, visitante.
  \item \textbf{Pregrado:} ninguna de las anteriores y algún estudiante llegó a 8.º semestre o más.
  \item \textbf{Pregrado nuevo:} la regla anterior excluía los programas recién creados, donde
  nadie ha llegado aún a 8.º. Se añadió una segunda señal, medida sobre las matrículas del
  programa: mediana de 12 o más créditos, mediana de valor pagado superior a \$500.000 y
  facultad asignada.
\end{itemize}

Quedaron \textbf{113 programas marcados como dudosos}; por defecto quedaban fuera. El consultor
revisó la propuesta y tomó \textbf{9 decisiones} que quedaron escritas, con su motivo, en
\cod{02c_filtro_validado.py}: confirmó dos pregrados nuevos que la regla había detectado pero
eran dudosos, incluyó los dos intercambios de pregrado (los únicos dudosos que entraron), excluyó
un «programa» de una sola persona y cuatro figuras de posgrado o visitantes. Aparte de esas 9,
decidió \textbf{no colapsar los programas pequeños}.

\begin{table}[!ht]
\centering
\caption{Universo final del estudio.}
\label{tab:universo}
\footnotesize
\begin{tabularx}{\columnwidth}{@{}L P{3.7cm}@{}}
\toprule
\cab{Universo final} & \cab{Valor} \\
\midrule
Programas & \textbf{60} de 576: 50 por la regla de pregrado, 8 pregrados nuevos, 2 intercambios \\
Filas & 448.489 (68,3\,\% del archivo) \\
Estudiantes & 18.563 \\
Matrículas (estudiante $\times$ programa $\times$ periodo) & 78.625 \\
Filas con matrícula simbólica ($\leq$ \$1.000) & 3.231 (0,72\,\%), que son 1.131 matrículas; se conservan marcadas \\
Programas con menos de 10 estudiantes & 14, con 66 personas; 2 tienen un solo estudiante \\
\bottomrule
\end{tabularx}
\end{table}

Un programa entró por la regla automática con \textbf{todas} sus matrículas simbólicas: «Sígueme
Tomasino -- Pregrado» (38 estudiantes, 45 matrículas). Quedan marcadas y, como las demás
simbólicas, deben excluirse del costeo.

\paso{4}{Seudonimizar de forma irreversible (Tarea 4)}

\begin{dorada}
\textbf{La idea.} Piense en una máquina de sellos que convierte cada cédula en un código. La
máquina se configura con una clave secreta elegida al azar entre $2^{256}$ ($\approx 10^{77}$)
posibilidades. Mientras existe la clave, la misma cédula produce siempre el mismo código, y así
se reconoce al mismo estudiante en todos los periodos. Cuando la clave se destruye, nadie
---tampoco quien la usó--- puede saber a qué cédula corresponde un código, ni calcular el código
de una cédula nueva.
\end{dorada}

\textbf{Por qué no un hash simple.} Una cédula colombiana tiene como mucho 10 dígitos: hay a lo
sumo $10^{10}$ posibilidades. Cualquiera puede calcular el SHA-256 de todas ellas en horas con un
equipo corriente y armar la tabla inversa. Con HMAC eso exige conocer la clave.

\textbf{Cómo se hizo}, dentro de \cod{03_anonimizar.py}:

\begin{enumerate}[leftmargin=*,itemsep=2pt,topsep=0pt]
  \item \textbf{Normalizar} el documento ---y, con la misma regla, el código de estudiante---:
  dejar solo dígitos y quitar los ceros a la izquierda (\cod{0012.345.678} $\to$ \cod{12345678};
  ejemplo ficticio). Esto fusionó a las 3 personas duplicadas: sin este paso habrían recibido
  dos códigos cada una y su historia habría salido partida.
  \item \textbf{Generar la llave}: \cod{secrets.token_bytes(32)}, el generador criptográfico del
  sistema operativo. 256 bits.
  \item \textbf{Calcular} \texttt{HMAC-SHA256(llave, etiqueta | documento)} y quedarse con los
  primeros caracteres hexadecimales (Tabla~\ref{tab:seudonimos}).

  Las \textbf{etiquetas} separan los dominios: si un mismo número apareciera como cédula y como
  código, sus seudónimos no se parecerían en nada, así que no se pueden cruzar las dos columnas.

  \cod{ID_PROGRAMA_EST} identifica una \textbf{vinculación}: la relación de una persona con un
  programa, que en el sistema académico es un código de estudiante. La etiqueta interna
  \cod{matricula} y el prefijo \cod{M} vienen de ese uso del sistema; no hay que confundirlos con
  la \emph{matrícula} de este documento (estudiante $\times$ programa $\times$ periodo).
  \item \textbf{Comprobar que no hay colisiones} (dos personas con el mismo código) sobre el
  archivo completo, antes de filtrar: 41.589 $\to$ 41.589 y 54.637 $\to$ 54.637. El script
  abortaba si había una. Por azar, la probabilidad era del orden de 0,0003\,\% para
  \cod{ID_EST} y 0,14\,\% para \cod{ID_PROGRAMA_EST}.
  \item \textbf{Destruir la llave}: \texttt{del llave}, y el proceso terminó al acabar el script. El
  código no la escribe en ningún sitio ---disco, log, variable de entorno--- y la auditoría no
  encontró rastros en ningún archivo. La única tabla de correspondencias cédula $\to$ código
  existió en la memoria del proceso y murió con él.
  \item \textbf{Prefijar} (hecho después, en el paso~10): \cod{E} para \cod{ID_EST}, \cod{M}
  para \cod{ID_PROGRAMA_EST}.
\end{enumerate}

\begin{table}[!ht]
\centering
\caption{Los dos seudónimos.}
\label{tab:seudonimos}
\footnotesize
\begin{tabularx}{\columnwidth}{@{}P{2.3cm} L L@{}}
\toprule
 & \cab{\cod{ID_EST}} & \cab{\cod{ID_PROGRAMA_EST}} \\
\midrule
Viene de & cédula normalizada & código de estudiante \\
Etiqueta & \cod{estudiante} & \cod{matricula} \\
Largo & 12 hex (48 bits) & 10 hex (40 bits) \\
Distintos en el archivo & 41.589 & 54.637 \\
Distintos en el entregable & 18.563 & 19.782 \\
Prefijo (paso~10) & \cod{E} & \cod{M} \\
\bottomrule
\end{tabularx}
\end{table}

El truncado a 12 y 10 caracteres \textbf{no debilita la irreversibilidad}, que depende de la
llave, no del largo del código. Solo afecta la probabilidad de colisión, que se comprobó.

\paso{5}{Suprimir y generalizar (Tarea 5)}

Sobre las 448.489 filas de pregrado se aplicaron los tratamientos de la
Tabla~\ref{tab:tratamientos}.

\begin{table*}[tp]
\centering
\caption{Supresiones y generalizaciones aplicadas a las filas de pregrado.}
\label{tab:tratamientos}
\footnotesize
\setlength{\tabcolsep}{4pt}
\begin{tabularx}{\textwidth}{@{}P{3.4cm} L P{2.1cm} P{4.6cm}@{}}
\toprule
\cab{Columna} & \cab{Tratamiento} & \cab{Afectados} & \cab{Por qué} \\
\midrule
6 identificadores & Suprimidos (cédula y código sobreviven como seudónimos) & todas las filas & Identifican directamente \\
19 de notas, promedios y fallas & Suprimidas, incluida la \cod{20221} & todas las filas & No sirven para costear; alto impacto \\
\cod{CAU} & Suprimida & --- & Constante: una sola sede \\
\cod{GRUPO} & Suprimida & --- & Singulariza (grupo $\times$ materia $\times$ periodo) y no interviene en el costeo \\
\cod{VAL_PAGADO} & Redondeado al millar más cercano (las mitades exactas, al millar par) & todas & Quita la huella de los últimos dígitos; error máximo \$500. Las simbólicas de \$100 quedan en \$0 \\
\cod{MATRICULA_SIMBOLICA} & Nueva bandera: valor original $\leq$ \$1.000 & 1.131 matrículas & Analítica: excluirlas del costeo \\
\cod{PER_INGRESO} & Ingresos anteriores a 2015 agrupados en \cod{<2015} & 1.132 matrículas & La cola antigua tiene pocas personas por periodo \\
\texttt{\#REIN} $\to$ \cod{N_REINTEGROS} & Valores de 4 o más (llegan a 14) agrupados en \cod{4+} & 1.117 matrículas & Cola larga y rara \\
\texttt{SER PILO PAGA} $\to$ \cod{SUBSIDIO_EXTERNO} & Verdadero/falso, sin nombrar el programa & 675 matrículas & Atributo socioeconómico de un grupo pequeño \\
\cod{Estado} $\to$ \cod{ESTADO} & \texttt{Sanción Disciplinaria} pasa a \cod{OTRO} & 36 matrículas & Dato disciplinario, raro y de alto impacto \\
\cod{N/A} $\to$ \cod{ES_NUEVO}; \cod{REINTEGRO} & Verdadero/falso & --- & Simplificación \\
\bottomrule
\end{tabularx}
\end{table*}

\textbf{Lo que no se generalizó, a propósito:} programa, periodo, nivel, facultad y las materias.
Son la sustancia del análisis. El plan proponía colapsar los programas pequeños; el consultor
decidió no hacerlo (decisión~6).

\paso{6}{Construir las tablas y barajar (Tarea 6)}

\cod{matriculas_anon} agrega el detalle por (\cod{ID_EST}, \cod{ID_PROGRAMA_EST},
\cod{PERIODO}) con las reglas de la Tabla~\ref{tab:agregacion}.

\begin{table}[!ht]
\centering
\caption{Reglas de agregación de \texttt{matriculas\_anon}.}
\label{tab:agregacion}
\footnotesize
\begin{tabularx}{\columnwidth}{@{}P{3.0cm} L@{}}
\toprule
\cab{Columna} & \cab{Regla} \\
\midrule
\cod{CREDITOS_INSCRITOS} & Suma de los créditos de las materias \\
\cod{N_MATERIAS} & Número de materias inscritas \\
\cod{N_MATERIAS_REPITE} & Materias que no se cursan por primera vez (\cod{VEZVISTA} distinto de \texttt{1-…}) \\
\cod{NIVEL_EST}, \cod{SUBSIDIO_EXTERNO}, \cod{REINTEGRO} & Máximo (da lo mismo que el primero, porque son constantes; se usó por prudencia) \\
Las demás (\cod{VAL_PAGADO}, programa, estado, etc.) & Primer valor \\
\bottomrule
\end{tabularx}
\end{table}

Tomar el «primer valor» solo es correcto si la columna es constante dentro de cada matrícula.
\textbf{Se verificó, no se supuso:} \cod{VAL_PAGADO} tiene un único valor por matrícula en las
78.625, y en la revisión del 2026-10-01 se comprobó lo mismo para las otras 12 columnas de origen
que se agregan (\cod{08_revision_documental.py} §2). La agregación no pierde información.

\cod{inscripciones_anon} conserva el grano original (una fila por materia) con 14 columnas.

\textbf{Barajado.} Las dos tablas se reordenaron al azar (semillas 20260903 y 20260904). Sin esto,
el orden de las filas seguiría el del archivo original, que viene ordenado de forma no aleatoria,
y la posición de una fila sería en sí misma una pista. Ver la limitación en la §\ref{sec:fuente}.

\textbf{Formato.} CSV separado por \cod{;}, UTF-8 con BOM, comprimido con gzip.

\paso{7}{Verificar el entregable (Checkpoint B)}

\cod{03b_verificar_entrega.py} hace \textbf{23 comprobaciones} y todas pasan
(Tabla~\ref{tab:checkpointb}). Eran 22: la del redondeo se rehízo el 2026-10-01 porque no leía
el entregable (§\ref{sec:incidentes}, R5).

\begin{table*}[tp]
\centering
\caption{Comprobaciones del Checkpoint B (\texttt{03b\_verificar\_entrega.py}).}
\label{tab:checkpointb}
\footnotesize
\setlength{\tabcolsep}{4pt}
\begin{tabularx}{\textwidth}{@{}P{2.4cm} L P{5.2cm}@{}}
\toprule
\cab{Grupo} & \cab{Qué se comprueba} & \cab{Resultado} \\
\midrule
Forma & Filas y columnas de cada tabla & 78.625 $\times$ 20 y 448.489 $\times$ 14 \\
Sin identificadores & Ninguna columna prohibida ni de notas; ningún valor de texto del entregable es igual a alguno de los 18.557 nombres distintos del universo & 0 nombres presentes \\
Dinero & Dentro de cada programa $\times$ periodo, los valores entregados son exactamente los originales redondeados; total original contra entregado & 497 de 497 grupos idénticos; \$394.563.422.461 $\to$ \$394.562.395.000 \\
Créditos & Original = detalle = agregado & 1.158.935 en los tres \\
Consistencia & Suma de \cod{N_MATERIAS} = filas del detalle; 11 periodos; 60 programas & Correcto \\
Integridad & Llave de matrícula única; inscripciones sin matrícula & Única; 0 huérfanas \\
Barajado & \cod{matriculas_anon} ya no está ordenada por periodo & Correcto \\
\bottomrule
\end{tabularx}
\end{table*}

La prueba del dinero merece una nota. Sin la llave no se puede emparejar cada matrícula entregada
con la suya en el original, así que se compara lo que sí es comparable: dentro de cada programa y
periodo, el conjunto de valores entregados tiene que ser idéntico al de los originales
redondeados. Lo es en los 497 grupos, y la prueba es sensible: con un solo valor alterado en
\$1.000, falla (\cod{08_revision_documental.py} §10). La diferencia total (0,00026\,\%) se
reporta, pero por sí sola no prueba nada.

\paso{8}{Medir el riesgo residual (Tarea 7)}

\cod{04_riesgo.py} mide sobre los entregables tres formas de reidentificación: combinaciones de
cuasi-identificadores (k-anonimato), unicidad del conjunto de materias y unicidad del valor
pagado. Resultados en la §\ref{sec:riesgo}.

\paso{9}{Diccionario y checksums (Tarea 8)}

\cod{05_diccionario.py} escribe el diccionario \textbf{leyendo los CSV reales} ---tipos, nulos,
rangos, valores frecuentes---, así esas columnas no pueden contradecir lo que se entrega. Origen y
tratamiento de cada columna salen de una tabla fija dentro del script, cuyas cifras se revisaron
el 2026-10-01. Escribe también \cod{CHECKSUMS.txt}.

\paso{10}{Auditoría final y corrección de los seudónimos (Tarea 9)}

\cod{06_auditoria_final.py} vuelve a verificar el original, barre los archivos de
\cod{anonimizacion/} ---salvo los intermedios en claro, que solo inventaría--- buscando
\textbf{formas} de correo, documento y celular, y busca rastros escritos de la llave. Encontró
tres problemas (§\ref{sec:incidentes}): documentos reales en la propia documentación, seudónimos
con forma de cédula y el auditor detectándose a sí mismo. Los tres se corrigieron.

La corrección de los seudónimos (\cod{07_prefijar_ids.py}) no podía hacerse regenerando: la llave
ya no existía y una nueva ejecución habría dado códigos distintos. Se \textbf{reetiquetó}:
anteponer \cod{E} o \cod{M} a cada código es una correspondencia uno a uno sobre los valores ya
emitidos, que conserva el vínculo entre periodos y entre las dos tablas (verificado: mismos
18.563 estudiantes, 19.782 vinculaciones, 0 inscripciones huérfanas).

\paso{11}{Revisión documental (2026-10-01)}

Al escribir este documento se recalcularon todas las cifras, se repitieron las verificaciones y se
revisó dónde hay copias en claro. Salieron un hallazgo grave (la semilla pública del barajado) y
cinco menores. Uno de ellos ---que una comprobación del Checkpoint B no leía el entregable--- lo
encontró una revisión independiente del borrador, que solo leyó documentos y scripts. Están en la §\ref{sec:incidentes} y,
con su evidencia, en la bitácora; \cod{08_revision_documental.py} los reproduce.

% ============================================================================
\section{Destino de las 52 columnas}\label{sec:columnas}
% ============================================================================

El destino de cada una de las 52 columnas del original ---tratamiento, columna resultante y tabla
en que queda--- está en el Anexo~\ref{app:columnas}. En resumen (Tabla~\ref{tab:cuenta}):

\begin{table}[!ht]
\centering
\caption{Cuenta de columnas.}
\label{tab:cuenta}
\footnotesize
\begin{tabularx}{\columnwidth}{@{}P{2.0cm} r L@{}}
\toprule
\cab{Movimiento} & \cab{N.º} & \cab{Cuáles} \\
\midrule
Salen & 31 & 6 identificadores, 19 de notas, \cod{CAU}, \cod{GRUPO} y 4 sustituidas por una versión recodificada \\
Siguen & 21 & Varias transformadas en su sitio \\
Se crean & 10 & 2 seudónimos, 4 versiones recodificadas, la bandera de matrícula simbólica y 3 conteos por matrícula \\
\midrule
Entregadas & 31 & 20 en \cod{matriculas_anon} y 14 en \cod{inscripciones_anon}; comparten 3 \\
\bottomrule
\end{tabularx}
\end{table}

\textbf{Cuenta.} Salen 31 columnas: 6 identificadores, 19 de notas, \cod{CAU}, \cod{GRUPO} y las 4
sustituidas por una versión recodificada (\cod{Estado}, \texttt{SER PILO PAGA}, \texttt{\#REIN} y
\cod{N/A}). Siguen 21, varias transformadas en su sitio. Se crean 10: los dos seudónimos, las
cuatro versiones recodificadas, la bandera de matrícula simbólica y tres conteos por matrícula.
Entre las dos tablas hay 31 columnas distintas ---la coincidencia con las 31 que salen es
casual---: 20 en \cod{matriculas_anon} y 14 en \cod{inscripciones_anon}, que comparten
\cod{ID_EST}, \cod{ID_PROGRAMA_EST} y \cod{PERIODO}.

% ============================================================================
\section{Qué se comprobó y cómo}\label{sec:comprobado}
% ============================================================================

La Tabla~\ref{tab:comprobado} reúne cada propiedad que se afirma del proceso, cómo se comprobó y
qué script lo hace.

\begin{table*}[tp]
\centering
\caption{Propiedades comprobadas y dónde se comprueban.}
\label{tab:comprobado}
\footnotesize
\setlength{\tabcolsep}{4pt}
\begin{tabularx}{\textwidth}{@{}P{3.5cm} L P{3.0cm} P{3.6cm}@{}}
\toprule
\cab{Propiedad} & \cab{Cómo se comprobó} & \cab{Resultado} & \cab{Dónde} \\
\midrule
El original nunca cambió & SHA-256 en cuatro momentos (inicio, tras extraer, auditoría, 2026-10-01) & Idéntico; permisos \cod{444} & \cod{06_auditoria_final.py} \\
La extracción es fiel & 1.300 celdas al azar contra el original; conteos de filas y de no nulos & Sin diferencias & \cod{01b_verificar.py} \\
No hay identificadores directos en los entregables & Columnas prohibidas; ningún valor igual a un nombre real; barrido de formas & 0 & \cod{03b_verificar_entrega.py}, \texttt{06\_…} \\
Ni la documentación ni los scripts contienen datos personales & Barrido de formas (correo, documento, celular) en \cod{anonimizacion/}; cruce de 18 archivos, incluidos los de la raíz, contra los nombres, documentos, correos y celulares reales & 0 y 0 & \cod{06_auditoria_final.py}; \cod{08_revision_documental.py} §7 \\
La llave no quedó escrita & Búsqueda de patrones de llave en scripts, logs y documentos & 0 rastros & \cod{06_auditoria_final.py} \\
Seudónimos sin colisiones & Distintos antes = distintos después, sobre el archivo completo & 0 colisiones & \cod{03_anonimizar.py} (aborta si las hay) \\
El análisis no se degradó & Valores por programa y periodo idénticos al original redondeado; créditos & 497 de 497 grupos; créditos idénticos & \cod{03b_verificar_entrega.py} \\
La agregación es exacta & Cada columna agregada es constante dentro de la matrícula & 0 excepciones en 13 columnas & \cod{08_revision_documental.py} §2 \\
Los entregables no cambian sin que se note & SHA-256 de cada archivo & Coinciden con \cod{CHECKSUMS.txt} & \cod{entrega/CHECKSUMS.txt} \\
\bottomrule
\end{tabularx}
\end{table*}

% ============================================================================
\section{Riesgo residual}\label{sec:riesgo}
% ============================================================================

Ninguna anonimización útil tiene riesgo cero. Esto es lo que queda, medido sobre los entregables
(informe completo en \cod{RIESGO_REIDENTIFICACION.md}).

\subsection{Combinaciones de atributos (k-anonimato)}

$k$ es cuántas matrículas comparten la misma combinación de atributos, contando la propia. Con
$k = 1$ la combinación señala a una sola persona. El plan fijó como referencia $k \geq 5$. Cada
fila de la Tabla~\ref{tab:kanon} suma atributos al conjunto de referencia por separado; no son
acumulativas.

\begin{table}[!ht]
\centering
\caption{Matrículas con combinación única ($k=1$) o casi única ($k<5$).}
\label{tab:kanon}
\footnotesize
\begin{tabularx}{\columnwidth}{@{}L r r@{}}
\toprule
\cab{Lo que sabe el atacante} & \cab{$k = 1$} & \cab{$k < 5$} \\
\midrule
Periodo y programa & 0,04\,\% & 0,18\,\% \\
\textbf{Periodo, programa y nivel (semestre)} --- conjunto de referencia & \textbf{0,45\,\%} & \textbf{2,91\,\%} \\
Referencia + periodo de ingreso & 6,88\,\% & 19,28\,\% \\
Referencia + valor pagado & 14,59\,\% & 33,72\,\% \\
\bottomrule
\end{tabularx}
\end{table}

\textbf{El criterio del plan (menos del 1\,\% bajo $k = 5$) no se cumple.} Era previsible: es la
consecuencia de conservar con su nombre los programas pequeños (decisión~6), y el consultor lo
aceptó en el Checkpoint A. El riesgo se concentra en licenciaturas en extinción y programas
pequeños.

\subsection{El conjunto de materias como huella}

Quien sabe qué materias cursó alguien puede encontrarlo, aunque no haya nombres
(Tabla~\ref{tab:materias}).

\begin{table}[!ht]
\centering
\caption{Unicidad del conjunto de materias.}
\label{tab:materias}
\footnotesize
\setlength{\tabcolsep}{4pt}
\begin{tabularx}{\columnwidth}{@{}L P{2.0cm} r@{}}
\toprule
\cab{Lo que sabe el atacante} & \cab{Se cuenta sobre} & \cab{Singulariza} \\
\midrule
3 materias sueltas, sin saber cuándo & estudiantes & 4,6\,\% \\
3 materias de un semestre concreto (un compañero de clase) & estudiantes en un periodo & 26,0\,\% \\
Todas las materias de un semestre & estudiantes en un periodo & 54,9\,\% \\
La historia completa de materias & estudiantes & 71,0\,\% \\
\bottomrule
\end{tabularx}
\end{table}

Este riesgo vive en \cod{inscripciones_anon} y \textbf{no se puede eliminar sin destruir el
detalle por materia}. Por eso \cod{matriculas_anon} es el entregable principal.

\subsection{El valor pagado}

El 5,0\,\% de las matrículas tiene un valor único dentro de su programa y periodo. Quien sepa
cuánto pagó una persona ---ella misma, su familia, quien maneje cartera--- puede encontrar su
matrícula. \cod{VAL_PAGADO} es a la vez la variable central del estudio y el cuasi-identificador
más fuerte: suprimirlo dejaría el dataset sin propósito.

\subsection{Escenarios}

\begin{table}[!ht]
\centering
\caption{Riesgo según lo que sabe quien intenta reidentificar.}
\label{tab:escenarios}
\footnotesize
\begin{tabularx}{\columnwidth}{@{}L P{3.4cm}@{}}
\toprule
\cab{Quién intenta reidentificar} & \cab{Riesgo} \\
\midrule
Alguien sin información previa & \textbf{Muy bajo}: no hay identificadores y la llave no existe \\
Alguien que conoce programa, periodo y semestre de una persona & Bajo en programas grandes, \textbf{alto en los pequeños} \\
Alguien que sabe cuánto pagó una persona & \alto\ \textbf{Alto} \\
Alguien que sabe qué materias cursó una persona & \alto\ \textbf{Muy alto} en \cod{inscripciones_anon} \\
Personal de la USTA con acceso a los sistemas fuente, sin la semilla & \alto\ \textbf{Alto}: cruzando columnas singulariza el 45,6\,\% de las matrículas (§\ref{sec:fuente}) \\
Quien tenga el archivo fuente \textbf{y} la semilla publicada & \alto\ \textbf{Total} en \cod{inscripciones_anon} y, a través de ella, en \cod{matriculas_anon} (§\ref{sec:fuente}) \\
\bottomrule
\end{tabularx}
\end{table}

\subsection{Recomendaciones}

\begin{enumerate}[leftmargin=*,itemsep=2pt,topsep=0pt]
  \item Compartir \textbf{\cod{matriculas_anon}} como producto principal.
  \item Compartir \textbf{\cod{inscripciones_anon}} solo si el análisis baja a nivel de materia, y
  con justificación escrita.
  \item Si los datos van a salir del equipo consultor, reconsiderar colapsar los 14 programas con
  menos de 10 estudiantes.
  \item Tratar ambos archivos como confidenciales: la anonimización reduce el riesgo, no lo anula.
  \item Antes de que salgan del equipo, \textbf{rebarajarlos} con una permutación que no se
  registre (§\ref{sec:pendientes}).
\end{enumerate}

% ============================================================================
\section{Límites: lo que esta anonimización no garantiza}\label{sec:limites}
% ============================================================================

\subsection{No hay garantía formal frente a cualquier atacante}

El método reduce y mide el riesgo; no ofrece garantías matemáticas como la privacidad
diferencial. Frente a alguien con suficiente información auxiliar sobre una persona concreta (sus
materias, su valor, su programa pequeño), el riesgo de la §\ref{sec:riesgo} es real.

\subsection{Frente a quien tenga el archivo fuente, el entregable no es anónimo}\label{sec:fuente}

La llave destruida protege frente a quien solo tiene los entregables. Pero quien tiene el
original ---la USTA, y este equipo mientras lo conserve--- puede volver a ligar los registros:

\begin{itemize}[leftmargin=*,itemsep=2pt,topsep=0pt]
  \item \textbf{Exactamente, con la semilla.} La semilla del barajado de \cod{inscripciones_anon}
  está escrita en \cod{03_anonimizar.py}. Con el original, el filtro de pregrado y esa semilla se
  rehace la permutación: se verificó que \textbf{448.489 de 448.489 filas} vuelven a su posición
  original. Cada \cod{ID_EST} queda ligado a su cédula y, como las dos tablas comparten
  \cod{ID_EST}, también \cod{matriculas_anon} cuando se tienen ambas.
  \item \textbf{Probablemente, sin la semilla.} El orden de \cod{matriculas_anon} no se puede
  rehacer así: antes de barajarla, la tabla estaba ordenada por los seudónimos, y ese orden solo
  se puede reconstruir con la llave. Pero cruzando siete columnas que el original también tiene
  ---periodo, programa, nivel, valor, créditos, número de materias y estado--- el 45,6\,\% de las
  matrículas es única, y una sola matrícula única liga toda la historia de esa persona.
\end{itemize}

\begin{alerta}
La semilla solo convierte en exacto y total algo que ya era probable. La conclusión práctica:
\textbf{para el poseedor del fuente esto es un dato seudonimizado}, y la forma de cerrarlo es no
conservar el fuente al terminar el proyecto. Rebarajar con una permutación no registrada elimina
la vía exacta; la probabilística no tiene remedio mientras exista el original.
\end{alerta}

\subsection{La llave existió en memoria}

\texttt{del llave} borra el nombre, no la memoria: Python no sobrescribe los bytes. La llave y la
tabla cédula $\to$ código vivieron en la memoria del proceso hasta que terminó el script. Lo
mitigan que el proceso fue único y breve, que macOS cifra la memoria de intercambio y que el disco
tiene FileVault. Lo que sí se garantiza y se verificó es que \textbf{nunca se escribió}.

\subsection{Los entregables no se pueden regenerar}

Es la otra cara de la irreversibilidad. Si se pierden o se dañan, una nueva ejecución dará otros
códigos, y cualquier análisis ya hecho sobre los actuales no se podrá cruzar con ella. Hay que
custodiar los archivos y verificar su checksum antes de usarlos. No se debe volver a ejecutar
\cod{03_anonimizar.py}, que generaría códigos nuevos, ni \cod{07_prefijar_ids.py}, que
reescribiría los archivos con otros checksums aunque el contenido fuera el mismo
(§\ref{sec:reejecutar}).

\subsection{La anonimización no corrige el dato}

Los problemas de calidad del archivo siguen ahí: \cod{PAGO} no informa nada; \cod{PER_INGRESO}
falta en el 43,6\,\% de las matrículas del entregable; hay 1 matrícula sin valor; y
\cod{N_REINTEGROS} no cuadra con \cod{REINTEGRO} en 8.148 matrículas (todo indica que el primero
es un conteo acumulado y el segundo el evento del periodo, sin confirmar con la USTA). Están
documentados en el diccionario.

\subsection{Esto no es un concepto jurídico}

La calificación de los entregables bajo la Ley 1581 ---en particular a la luz de la
§\ref{sec:fuente}--- la debe hacer la USTA.

% ============================================================================
\section{Custodia de los datos en claro, destrucción y decisiones pendientes}\label{sec:custodia}
% ============================================================================

\subsection{Dónde hay datos personales en claro hoy}

La Tabla~\ref{tab:inventario} es el inventario de copias en claro a la fecha de la revisión.

\begin{table*}[tp]
\centering
\caption{Inventario de datos personales en claro (2026-10-01).}
\label{tab:inventario}
\footnotesize
\setlength{\tabcolsep}{4pt}
\begin{tabularx}{\textwidth}{@{}P{4.6cm} L Q{1.45cm} P{5.2cm}@{}}
\toprule
\cab{Ubicación} & \cab{Contenido} & \cab{Tamaño} & \cab{Destino previsto} \\
\midrule
\cod{OFERTA_ACADEMICA_20211_20261.xlsx} (raíz del proyecto) & El original de la USTA, completo & 142,8 MB & Conservar como evidencia hasta el cierre; después devolver o eliminar según indique la USTA (P5) \\
\cod{anonimizacion/trabajo/OFERTA_COPIA.xlsx} & Copia bit a bit del original & 142,8 MB & Destruir (P2) \\
\cod{anonimizacion/_INSEGURO/oferta_raw.parquet} & Las 656.737 filas sin anonimizar & 19,7 MB & Destruir (P2) \\
\bottomrule
\end{tabularx}
\end{table*}

\cod{_INSEGURO/log_01_extraer.txt} está en la carpeta de intermedios pero solo contiene nombres de
columna y conteos. Los entregables, la documentación y los scripts no contienen datos personales
(auditoría y cruce contra los valores reales, §\ref{sec:comprobado}).

\subsection{Hallazgo de la revisión del 2026-10-01 que requiere una decisión}\label{sec:hallazgos}

\begin{alerta}
\textbf{\color{ustaRojo}R1 $\cdot$ La semilla pública permite deshacer el barajado} a quien tenga
el original (§\ref{sec:fuente}). La vía exacta se cierra rebarajando con una permutación que no
se registre (P1); la probabilística, no conservando el fuente al cierre (P5).
\end{alerta}

\subsection{Decisiones pendientes}\label{sec:pendientes}

La Tabla~\ref{tab:pendientes} lista las cinco decisiones abiertas. Ninguna de estas acciones se
ejecutó en la revisión: todas borran, reescriben o comprometen datos, y le corresponden a quien
indica la columna «Decide».

\begin{table*}[tp]
\centering
\caption{Decisiones pendientes.}
\label{tab:pendientes}
\footnotesize
\setlength{\tabcolsep}{4pt}
\begin{tabularx}{\textwidth}{@{}l P{2.7cm} P{1.9cm} L P{4.4cm}@{}}
\toprule
\cab{\#} & \cab{Decisión} & \cab{Decide} & \cab{Opciones} & \cab{Recomendación} \\
\midrule
P1 & \textbf{Rebarajar los entregables} con una permutación aleatoria que no se registre & Consultor & Hacerlo ya, o justo antes de que salgan del equipo & Antes de cualquier envío fuera del equipo. No cambia contenido ni seudónimos, solo el orden y los checksums \\
P2 & \textbf{Destruir los intermedios en claro} & Consultor & Ya, o al cierre del proyecto & Ya: si hicieran falta, se regeneran desde el original en un minuto \\
P3 & \textbf{Autorización institucional} del tratamiento & USTA, a solicitud del consultor & --- & Pedirla por escrito \\
P4 & \textbf{Aceptación escrita del riesgo} residual ($k < 5$ en el 2,91\,\%) & USTA & Someter el informe de riesgo o no & Someterlo, junto con la §\ref{sec:fuente}. El consultor ya lo aceptó en el Checkpoint A; falta la USTA \\
P5 & \textbf{El original al cierre} & USTA y consultor & Devolverlo o eliminarlo & Es lo único que mantiene abierta la vía de la §\ref{sec:fuente} \\
\bottomrule
\end{tabularx}
\end{table*}

\subsection{Procedimiento de destrucción al cierre}

Desde la raíz del proyecto:

\begin{codigo}[]
# 1. Auditar antes: debe decir APROBADA
python3 anonimizacion/scripts/06_auditoria_final.py
# 2. Borrar los intermedios
rm -rf anonimizacion/_INSEGURO anonimizacion/trabajo
\end{codigo}

\begin{enumerate}[leftmargin=*,itemsep=2pt,topsep=0pt,start=3]
  \item Decidir el destino del original (P5).
  \item Registrar en \cod{EVIDENCIA.md} la fecha y lo destruido.
\end{enumerate}

Con FileVault activo no hace falta sobrescribir el disco local.

% ============================================================================
\section{Incidentes, desviaciones y correcciones}\label{sec:incidentes}
% ============================================================================

Registrar los errores es parte de la evidencia: muestra qué se comprobó de verdad.

\subsection{Desviaciones respecto al plan}\label{sec:desviaciones}

La Tabla~\ref{tab:desviaciones} compara lo planeado con lo ejecutado.

\begin{table*}[tp]
\centering
\caption{Desviaciones respecto al plan aprobado.}
\label{tab:desviaciones}
\footnotesize
\setlength{\tabcolsep}{4pt}
\begin{tabularx}{\textwidth}{@{}L L L@{}}
\toprule
\cab{Plan} & \cab{Ejecutado} & \cab{Motivo} \\
\midrule
Dar tipo a los datos al extraer & Extraer texto crudo y dar tipo en \cod{03} & Verificación celda a celda exacta \\
«21 columnas de notas» & 19 & Error de conteo del plan \\
La columna \cod{20221} está vacía o es espuria & Tiene 283.443 calificaciones; se suprime como nota & Descubierto al perfilar \\
Scripts \cod{02_anonimizar} y \cod{03_tablas} & Un solo \cod{03_anonimizar.py} & La llave efímera exige un proceso \\
41.592 cédulas $\to$ 41.592 códigos & 41.589 & 3 personas duplicadas por ceros a la izquierda \\
\texttt{\#REIN} mayor que 3 $\to$ \cod{3+} & 4 o más $\to$ \cod{4+} & Mismo agrupamiento con una etiqueta exacta \\
Colapsar programas con pocos estudiantes & No se colapsaron & Decisión del consultor \\
Colapsar \texttt{Deserción Precoz} & No aparece en pregrado; solo se colapsó la sanción & Dato \\
Menos del 1\,\% bajo $k = 5$ & 2,91\,\% & Consecuencia aceptada de la decisión anterior \\
Si $k < 5$ supera el 1\,\%, iterar sobre la Tarea 5 con generalización más agresiva & No se iteró & Se aceptó el riesgo en lugar de colapsar programas \\
Revisión visual de 50 filas por el consultor (Checkpoint B) & \cod{03b} muestra 5 matrículas al azar & Simplificación no registrada en su momento \\
Seudónimos hexadecimales & Prefijados \texttt{E…} / \texttt{M…} & Incidente 2 \\
Destruir los intermedios en la Tarea 9 & Conservados & Decisión del consultor \\
\bottomrule
\end{tabularx}
\end{table*}

\subsection{Incidentes}

La Tabla~\ref{tab:incidentes} recoge los incidentes de la ejecución (1 a 3) y los hallazgos de la
revisión del 2026-10-01 (R1 a R6).

\begin{table*}[tp]
\centering
\caption{Incidentes y hallazgos, con su gravedad y corrección.}
\label{tab:incidentes}
\footnotesize
\setlength{\tabcolsep}{4pt}
\begin{tabularx}{\textwidth}{@{}P{1.7cm} L P{3.6cm} L@{}}
\toprule
\cab{\#} & \cab{Qué pasó} & \cab{Gravedad} & \cab{Corrección} \\
\midrule
Umbral no garantizado & La verificación exigía una diferencia total por redondeo menor al 0,0001\,\%, que el redondeo no garantiza & Bajo & Se pasó a comprobar la cota de \$500\ldots\ con una prueba que no leía el entregable (ver R5) \\
1 & Al documentar los duplicados se copiaron \textbf{7 documentos reales} (de 4 personas) en la bitácora y en un comentario de código & \alto\ Alto & Redactados; el ejemplo pasó a ser ficticio; la auditoría vigila las formas desde entonces \\
2 & 185 \cod{ID_PROGRAMA_EST} salieron todo-dígitos por azar y parecían cédulas o celulares & Bajo, falso positivo & Prefijos \cod{E} / \cod{M} sin regenerar \\
3 & La auditoría se detectaba a sí misma & Bajo & Se excluye del barrido; patrones con límites estrictos \\
R1 & La semilla pública permite deshacer el barajado a quien tenga el fuente & \alto\ \textbf{Alto} frente al poseedor del fuente & Pendiente: P1 y P5 \\
R2 & La auditoría documentada fallaba por un literal en la descripción del Incidente 3 & Bajo & Frase reescrita; la auditoría vuelve a aprobar \\
R3 & El informe de riesgo afirmaba sin medir que 3 materias singularizan a la mayoría & Bajo & Medido (4,6\,\% / 26,0\,\%); informe regenerado \\
R4 & Imprecisiones de cifras, unidades, enlaces y comentarios en bitácora, diccionario y scripts & Menor & Corregidas; los CSV no cambiaron \\
R5 & La comprobación «error $\leq$ \$500 por matrícula» comparaba el original consigo mismo y no leía el entregable & Bajo & Rehecha por programa $\times$ periodo: 497 de 497 grupos; se probó que detecta un solo valor alterado \\
R6 & El auditor perdía un documento seguido de punto final, solo busca formas y solo recorre \cod{anonimizacion/}; varias cifras de la revisión no tenían script & Bajo & Patrón corregido; cruce contra valores reales y cifras en \cod{08_revision_documental.py} \\
\bottomrule
\end{tabularx}
\end{table*}

\begin{nota}
\textbf{\color{ustaAzul}Lecciones.} La del Incidente~1 se convirtió en regla verificada:
\textbf{no se escriben identificadores reales en documentación ni en código, ni siquiera como
ejemplo}. La auditoría vigila las formas en cada corrida, y \cod{08_revision_documental.py}
cruza contra los valores reales. Y la de R5: una comprobación hay que probarla haciéndola fallar
a propósito.
\end{nota}

% ============================================================================
\section{Cómo verificar y qué se puede reproducir}\label{sec:verificar}
% ============================================================================

\subsection{Verificar}

Desde la raíz del proyecto se ejecutan los comandos de la Figura~\ref{fig:comandos}. El primero
debe decir \texttt{original INTACTO}; el segundo, \texttt{OK} en los tres archivos; el tercero,
\texttt{AUDITORÍA: APROBADA}. El cuarto repite la revisión completa del 2026-10-01 (unos dos
minutos; lee el intermedio en claro e imprime solo conteos).

\begin{figure*}[tp]
\begin{codigo}[basicstyle=\ttfamily\footnotesize]
# 1. El original no cambió
shasum -a 256 OFERTA_ACADEMICA_20211_20261.xlsx \
  | grep -q f447e0d24e716b09a9abea1655dec6a0c2e143bd418fd93115e826c5a0eaa9d9 \
  && echo "original INTACTO" || echo "ALERTA: el original cambió"

# 2. Los entregables coinciden con sus checksums
(cd anonimizacion/entrega && grep -v '^#' CHECKSUMS.txt | grep . | shasum -a 256 -c)

# 3. Auditoría final
python3 anonimizacion/scripts/06_auditoria_final.py

# 4. Revisión documental completa (solo imprime conteos)
python3 anonimizacion/scripts/08_revision_documental.py
\end{codigo}
\caption{Comandos de verificación, desde la raíz del proyecto.}
\label{fig:comandos}
\end{figure*}

\subsection{Qué se puede volver a ejecutar}\label{sec:reejecutar}

La Tabla~\ref{tab:reejecutar} dice, para cada script, si se puede repetir, si necesita datos en
claro y qué escribe.

\begin{table*}[tp]
\centering
\caption{Scripts: re-ejecución, datos que leen y archivos que escriben.}
\label{tab:reejecutar}
\footnotesize
\setlength{\tabcolsep}{4pt}
\begin{tabularx}{\textwidth}{@{}P{4.1cm} P{2.3cm} L L@{}}
\toprule
\cab{Script} & \cab{¿Se puede re-ejecutar?} & \cab{Necesita datos en claro} & \cab{Escribe} \\
\midrule
\cod{01_extraer.py} & Sí & Sí, la copia & \cod{_INSEGURO/oferta_raw.parquet} \\
\cod{01b_verificar.py} & Sí & Sí, el parquet y el original & nada \\
\cod{02_catalogo_programas.py} & Sí, determinista & Sí & \cod{catalogo_programas.csv} \\
\cod{02c_filtro_validado.py} & Sí, determinista & No & \cod{filtro_pregrado.csv} \\
\rowcolor{ustaRojo!7}\textbf{\cod{03_anonimizar.py}} & \nose & Sí & \textbf{Sobrescribe los entregables con códigos nuevos} \\
\cod{03b_verificar_entrega.py} & Sí & Sí & nada \\
\cod{04_riesgo.py} & Sí, determinista & No & \cod{RIESGO_REIDENTIFICACION.md} \\
\cod{05_diccionario.py} & Sí & No & \cod{DICCIONARIO_DATOS.md}, \cod{CHECKSUMS.txt} \\
\cod{06_auditoria_final.py} & Sí & Lee el original solo para calcular su huella; de los intermedios solo lee los \cod{.txt} y \cod{.md} al buscar rastros de la llave & nada \\
\rowcolor{ustaRojo!7}\textbf{\cod{07_prefijar_ids.py}} & \nose & No & Reescribe los CSV: mismo contenido, pero gzip guarda la hora y \textbf{cambian los checksums} \\
\cod{08_revision_documental.py} & Sí & Sí; imprime solo conteos & nada \\
\bottomrule
\end{tabularx}
\end{table*}

Si algún día hay que repetir la anonimización completa ---por ejemplo, con un archivo nuevo---
hay que tratarla como un entregable distinto, con códigos que no se cruzan con los actuales, y
quitar la semilla fija del barajado (§\ref{sec:fuente}).

% ============================================================================
% ANEXOS
% ============================================================================
\clearpage
\appendix
\onecolumn
\titleformat{\section}{\normalfont\large\bfseries\raggedright\color{ustaAzul}}{Anexo \thesection.}{0.5em}{}[{\color{ustaDorado}\titlerule[0.8pt]}]
\addtocontents{toc}{\protect\medskip\protect\noindent{\protect\bfseries\protect\color{ustaAzul}Anexos}\protect\par}

% ============================================================================
\section{Destino de las 52 columnas}\label{app:columnas}
% ============================================================================

\textbf{M} = \cod{matriculas_anon} \quad$\cdot$\quad \textbf{I} = \cod{inscripciones_anon}.

{\small
\begin{xltabular}{\textwidth}{@{}l P{4.6cm} L P{4.1cm} l@{}}
\caption{Destino de cada columna del archivo original.}\label{tab:columnas52}\\
\toprule
\cab{Col} & \cab{Columna original} & \cab{Tratamiento} & \cab{Columna resultante} & \cab{Tabla} \\
\midrule
\endfirsthead
\multicolumn{5}{l}{\textit{(continuación de la Tabla~\ref{tab:columnas52})}}\\
\toprule
\cab{Col} & \cab{Columna original} & \cab{Tratamiento} & \cab{Columna resultante} & \cab{Tabla} \\
\midrule
\endhead
\midrule
\multicolumn{5}{r}{\textit{(continúa en la página siguiente)}}\\
\endfoot
\bottomrule
\endlastfoot
A & \cod{CAU} & Suprimida (constante) & --- & --- \\
B & \cod{FACULTAD} & Sin cambios & \cod{FACULTAD} & M \\
C & \cod{PERIODO} & Sin cambios & \cod{PERIODO} & M, I \\
D & \cod{CODIGO} & \textbf{Seudonimizada} & \cod{ID_PROGRAMA_EST} & M, I \\
E & \cod{COD_UNIDAD} & Sin cambios & \cod{COD_UNIDAD} & M \\
F & \texttt{PROGRAMA ACADEMICO} & Filtrada a 60 programas; renombrada & \cod{PROGRAMA_ACADEMICO} & M \\
G & \texttt{APELLIDOS Y NOMBRES} & \textbf{Suprimida} & --- & --- \\
H & \cod{IDENTIFICACION} & \textbf{Seudonimizada} & \cod{ID_EST} & M, I \\
I & \cod{PER_INGRESO} & Generalizada en su sitio (\cod{<2015}) & \cod{PER_INGRESO} & M \\
J & \cod{DIR_EMAIL} & \textbf{Suprimida} & --- & --- \\
K & \cod{CELULAR} & \textbf{Suprimida} & --- & --- \\
L & \cod{DIR_EMAIL_PER} & \textbf{Suprimida} & --- & --- \\
M & \cod{PENSUM_ESTUDIANTE} & Sin cambios & \cod{PENSUM_ESTUDIANTE} & M \\
N & \cod{MAT_PENSUM} & Sin cambios & \cod{MAT_PENSUM} & I \\
O & \cod{MAT_PENSUM_1} & Sin cambios & \cod{MAT_PENSUM_1} & I \\
P & \cod{Estado} & Sustituida por una versión recodificada (sanción $\to$ \cod{OTRO}) & \cod{ESTADO} & M \\
Q & \cod{C_MAT} & Sin cambios & \cod{C_MAT} & I \\
R & \cod{PENSUM_OFERTA} & Sin cambios & \cod{PENSUM_OFERTA} & I \\
S & \cod{NOM_MATERIA} & Sin cambios & \cod{NOM_MATERIA} & I \\
T & \cod{GRUPO} & Suprimida (cuasi-identificador) & --- & --- \\
U & \cod{COD_UNIDAD_1} & Renombrada & \cod{COD_UNIDAD_OFERTA} & I \\
V & \texttt{programa que oferta la materia} & Renombrada & \cod{PROGRAMA_OFERTA} & I \\
W & \cod{FAC_OFERTA} & Sin cambios & \cod{FAC_OFERTA} & I \\
X & \cod{NIVEL_EST} & Máximo por matrícula & \cod{NIVEL_EST} & M \\
Y & \cod{NIVEL_MAT} & Sin cambios & \cod{NIVEL_MAT} & I \\
Z & \cod{CREDITOS} & Sin cambios; sumada por matrícula & \cod{CREDITOS} $\cdot$ \cod{CREDITOS_INSCRITOS} & I $\cdot$ M \\
AA & \cod{VAL_PAGADO} & Redondeada al millar en su sitio; deriva una bandera & \cod{VAL_PAGADO} $\cdot$ \cod{MATRICULA_SIMBOLICA} & M \\
AB & \cod{PAGO} & Sin cambios & \cod{PAGO} & M \\
AC & \cod{PROM_SEMESTRE} & Suprimida (nota) & --- & --- \\
AD & \cod{PRO_ACUMULADO} & Suprimida (nota) & --- & --- \\
AE & \cod{20221} & Suprimida (contiene notas) & --- & --- \\
AF & \cod{VEZVISTA} & Sin cambios; contada por matrícula & \cod{VEZVISTA} $\cdot$ \cod{N_MATERIAS_REPITE} & I $\cdot$ M \\
AG & \texttt{SER PILO PAGA} & Sustituida por una versión recodificada (verdadero/falso) & \cod{SUBSIDIO_EXTERNO} & M \\
AH & \cod{REINTEGRO} & Recodificada en su sitio (verdadero/falso) & \cod{REINTEGRO} & M \\
AI & \texttt{\#REIN} & Sustituida por una versión recodificada (topada en \cod{4+}) & \cod{N_REINTEGROS} & M \\
AJ & \cod{DEFINITIVA} & Suprimida (nota) & --- & --- \\
AK & \cod{N/A} (antiguo o nuevo) & Sustituida por una versión recodificada (verdadero/falso) & \cod{ES_NUEVO} & M \\
AL--AZ & \texttt{\%1}--\texttt{\%5}, \cod{Nota1}--\cod{Nota5}, \cod{Fallas1}--\cod{Fallas5} & Suprimidas (15 de notas) & --- & --- \\
--- & (nueva) & Conteo de materias por matrícula & \cod{N_MATERIAS} & M \\
\end{xltabular}
}

% ============================================================================
\section{Glosario}\label{app:glosario}
% ============================================================================

{\small
\begin{xltabular}{\textwidth}{@{}P{4.2cm} L@{}}
\caption{Términos usados en este documento.}\label{tab:glosario}\\
\toprule
\cab{Término} & \cab{Significado aquí} \\
\midrule
\endfirsthead
\multicolumn{2}{l}{\textit{(continuación de la Tabla~\ref{tab:glosario})}}\\
\toprule
\cab{Término} & \cab{Significado aquí} \\
\midrule
\endhead
\midrule
\multicolumn{2}{r}{\textit{(continúa en la página siguiente)}}\\
\endfoot
\bottomrule
\endlastfoot
\textbf{Anonimizar} & Lograr que nadie pueda ligar un registro con una persona por medios razonables \\
\textbf{Seudonimizar} & Reemplazar los identificadores por códigos; si el vínculo puede rehacerse, el dato sigue siendo personal \\
\textbf{Responsable / encargado} & Quien decide sobre el tratamiento de los datos (la USTA) y quien los trata por cuenta de él (el consultor) \\
\textbf{HMAC-SHA256} & Función que mezcla un dato con una llave secreta y da un código; sin la llave no se puede invertir ni recalcular \\
\textbf{Llave efímera} & Llave que solo existe durante la ejecución y nunca se guarda \\
\textbf{Hexadecimal, bit} & Un carácter hexadecimal (0--9, a--f) codifica 4 bits; 12 caracteres son 48 bits \\
\textbf{Colisión} & Que dos valores distintos reciban el mismo código \\
\textbf{Semilla} & Número que fija una secuencia «aleatoria»: con la misma semilla sale siempre la misma \\
\textbf{Permutación} & Un reordenamiento de las filas \\
\textbf{Barajar} & Reordenar las filas al azar para que la posición no sea una pista \\
\textbf{Cuasi-identificador} & Atributo que no identifica solo pero sí combinado (programa + periodo + semestre) \\
\textbf{k-anonimato} & Cuántos registros comparten una combinación de cuasi-identificadores, contando el propio; $k = 1$ es único \\
\textbf{Unicidad de trayectoria} & Que el conjunto de materias de alguien no lo tenga nadie más \\
\textbf{Privacidad diferencial} & Técnica que añade ruido calculado para dar una garantía matemática de privacidad; aquí no se usó \\
\textbf{Generalizar} & Hacer menos preciso un dato (redondear, agrupar) \\
\textbf{Topar} (\emph{top-coding}) & Agrupar los valores extremos en una categoría abierta (\cod{4+}) \\
\textbf{Checksum SHA-256} & Huella de 64 caracteres de un archivo; si cambia un solo bit, cambia la huella \\
\textbf{Intermedio en claro} & Archivo de trabajo que todavía tiene datos personales sin anonimizar \\
\textbf{Parquet} & Formato de archivo por columnas, compacto, para datos tabulares \\
\textbf{Lectura en streaming} & Leer un archivo por partes, sin cargarlo entero en memoria \\
\textbf{BOM} & Marca invisible al inicio de un archivo UTF-8 que le indica a Excel la codificación \\
\textbf{Matrícula} & Aquí: un estudiante en un programa en un periodo. Es el grano en que se paga \\
\textbf{Vinculación} & La relación de una persona con un programa (un código de estudiante). Quien cursa dos programas tiene dos \\
\textbf{Inscripción} & Una materia inscrita en una matrícula. Es el grano del archivo original \\
\textbf{Matrícula simbólica} & Valor pagado de \$1.000 o menos (intercambios, casos sin cobro ordinario) \\
\textbf{Nivel} & Semestre que cursa el estudiante (\cod{NIVEL_EST}) o al que pertenece la materia (\cod{NIVEL_MAT}) \\
\textbf{Pensum} & Plan de estudios de un programa en una versión dada \\
\textbf{Tarea, checkpoint} & Las 9 tareas del plan y sus 3 puntos de control con aprobación del consultor \\
\end{xltabular}
}

% ============================================================================
\section{Mapa de documentos y scripts}\label{app:mapa}
% ============================================================================

Rutas relativas a \cod{anonimizacion/}, salvo las que empiezan por \cod{../} (raíz del proyecto).

{\small
\begin{xltabular}{\textwidth}{@{}P{5.6cm} L P{4.4cm}@{}}
\caption{Documentos del proceso y cuándo leerlos.}\label{tab:mapa}\\
\toprule
\cab{Archivo} & \cab{Qué es} & \cab{Leer cuando\ldots} \\
\midrule
\endfirsthead
\multicolumn{3}{l}{\textit{(continuación de la Tabla~\ref{tab:mapa})}}\\
\toprule
\cab{Archivo} & \cab{Qué es} & \cab{Leer cuando\ldots} \\
\midrule
\endhead
\midrule
\multicolumn{3}{r}{\textit{(continúa en la página siguiente)}}\\
\endfoot
\bottomrule
\endlastfoot
\cod{METODOLOGIA_ANONIMIZACION.md} & La versión en Markdown de este documento & Siempre, primero \\
\cod{EVIDENCIA.md} & Bitácora cronológica con checksums, decisiones, incidentes y la revisión del 2026-10-01 & Se necesita la prueba de algo \\
\cod{RIESGO_REIDENTIFICACION.md} & Informe de riesgo residual, generado por \cod{04_riesgo.py} & Se va a compartir un entregable \\
\cod{entrega/DICCIONARIO_DATOS.md} & Columnas, tipos, rangos y advertencias de uso & Se van a analizar los datos \\
\cod{entrega/CHECKSUMS.txt} & Huellas de los entregables & Antes de usar o enviar un entregable \\
\cod{../PLAN_ANONIMIZACION.md} & El plan aprobado (2026-09-02), con sus 9 tareas & Interesa lo que se planeó (las diferencias con lo ejecutado están en la §\ref{sec:desviaciones}) \\
\cod{catalogo_programas.csv} & Los 576 programas clasificados (solo agregados) & Se cuestiona el filtro de pregrado \\
\cod{filtro_pregrado.csv} & Los 576 programas con la marca \cod{incluir_en_estudio} (60 verdaderas) y el motivo de cada decisión & Ídem \\
\cod{_INSEGURO/LEEME_ANTES_DE_TOCAR.md} & Advertencias sobre los intermedios en claro & Antes de abrir esa carpeta \\
\cod{../REGISTRO_PROYECTO.md} & Registro general del proyecto; su §8 tiene el chequeo de salud & Al retomar el proyecto \\
\end{xltabular}
}

Los scripts están en \cod{scripts/}, numerados en el orden en que se ejecutaron; la
Tabla~\ref{tab:reejecutar} dice qué hace cada uno y si se puede repetir.

% ============================================================================
\section{Cifras clave}\label{app:cifras}
% ============================================================================

{\small
\begin{xltabular}{\textwidth}{@{}L P{5.2cm} P{4.6cm}@{}}
\caption{Cifras clave y el script que las produce.}\label{tab:cifras}\\
\toprule
\cab{Cifra} & \cab{Valor} & \cab{Fuente} \\
\midrule
\endfirsthead
\multicolumn{3}{l}{\textit{(continuación de la Tabla~\ref{tab:cifras})}}\\
\toprule
\cab{Cifra} & \cab{Valor} & \cab{Fuente} \\
\midrule
\endhead
\midrule
\multicolumn{3}{r}{\textit{(continúa en la página siguiente)}}\\
\endfoot
\bottomrule
\endlastfoot
Filas y columnas del original & 656.737 $\times$ 52 & perfilado; \cod{01b_verificar.py} \\
Personas en el original & 41.589 (41.592 antes de normalizar) & \cod{03_anonimizar.py} \\
Programas en el original & 576 & \cod{02_catalogo_programas.py} \\
Programas del estudio & 60 & \cod{02c_filtro_validado.py} \\
Filas del estudio & 448.489 (68,3\,\%) & \cod{03_anonimizar.py} \\
Estudiantes del estudio & 18.563 & \cod{03b_verificar_entrega.py} \\
Vinculaciones a programa & 19.782 & \cod{07_prefijar_ids.py} \\
Matrículas & 78.625 & \cod{03b_verificar_entrega.py} \\
Valor total matriculado & \$394.563.422.461 (original) $\cdot$ \$394.562.395.000 (entregado) & \cod{03b_verificar_entrega.py} \\
Grupos programa $\times$ periodo con valores idénticos al original redondeado & 497 de 497 & \cod{03b_verificar_entrega.py} \\
Error máximo por redondeo & \$500 por matrícula & \cod{03b_verificar_entrega.py} \\
Créditos inscritos & 1.158.935 & \cod{03b_verificar_entrega.py} \\
Matrículas simbólicas & 1.131 (1.127 de \$100 y 4 de \$0) & \cod{08_revision_documental.py} §3 \\
Columnas suprimidas & 31 de 52 & \cod{03_anonimizar.py} \\
$k < 5$ (conjunto de referencia) & 2,91\,\% de las matrículas & \cod{04_riesgo.py} \\
Estudiantes con historia de materias única & 71,0\,\% & \cod{04_riesgo.py} \\
Filas re-ligables con fuente + semilla & 448.489 de 448.489 & \cod{08_revision_documental.py} §5 \\
Matrículas únicas para el poseedor del fuente, sin semilla & 45,6\,\% & \cod{08_revision_documental.py} §6 \\
\end{xltabular}
}

\vspace{1cm}
\begin{flushright}
{\color{ustaAzul}\bfseries Consultoría USTA 2026-II --- Universidad Santo Tomás}\\
Responsable: Javier Mauricio Sierra\\
1 de octubre de 2026
\end{flushright}

% ============================================================================
\end{document}
